% SEGMENT OLP-0638-S01
% Part: history
% Chapter: set-theory
% Section: cantor-plane

\documentclass[../../../include/open-logic-section]{subfiles}

\begin{document}

\olfileid{his}{set}{cantorplane}
\olsection{கோட்டையும் தளத்தையும் பற்றிய காண்டோரின் கருத்து}

ஷ்ரோடர்--பெர்ன்ஸ்டைன் தேற்றத்தின் நிறுவலைச் சூழ்ந்த சில
நிகழ்வுகள், \olref[his][set][pathology]{sec} பகுதியில் நாம்
விவாதித்த வரலாற்றுடன் தொடர்புடையவை. 1877-இல் ஒரு சதுரத்தில்
அதன் ஒரு பக்கத்தில் உள்ள அதே எண்ணிக்கையிலான புள்ளிகள் உள்ளன
என்று காண்டோர் நிறுவியதை நினைவுகூருங்கள். இங்கே அவரது
(முதல் முயற்சியான) நிறுவலைக் காண்போம்.

$\unitline$ என்பது அலகுக் கோடு, அதாவது $[0,1]$-இல் உள்ள
புள்ளிகளின் கணம் என்று கொள்வோம். $\unitsquare$ என்பது அலகுச்
சதுரம், அதாவது $\unitline
\times \unitline$-இல் உள்ள புள்ளிகளின் கணம் என்று கொள்வோம்.
இக்குறியீட்டில், காண்டோர்
$\cardeq{\unitline}{\unitsquare}$ என்று நிறுவினார்.
டெடெகிண்டிற்கு அவர் எழுதிய குறிப்பில் பின்வரும் வாதம்
அடிப்படையில் இருந்தது.

% SEGMENT OLP-0638-S02
\begin{thm}\ollabel{thm:cantorplane}
$\cardeq{\unitline}{\unitsquare}$
\end{thm}

\begin{proof}[நிறுவல்: முதல் பகுதி.]
$a, b \in \unitline$ என்று நிலைப்படுத்துவோம். அவற்றை
இரும எண் குறியீட்டில் எழுதுவோம். அப்போது $0$-களும்
$1$-களும் அடங்கிய முடிவிலாத் தொடர்களான
$a_1$, $a_2$, \dots, மற்றும் $b_1$, $b_2$, \dots, கிடைக்கும்:
\begin{align*}
a &= 0.a_1a_2a_3a_4\dots\\
b &= 0.b_1b_2b_3b_4\dots
\intertext{இப்போது $f \colon \unitsquare \to \unitline$ என்ற சார்பைக் கீழ்வருமாறு வரையறுப்போம்:}
f(a, b) & = 0.a_1b_1a_2b_2a_3b_3a_4b_4\dots
\end{align*}
இப்போது $f$ ஓர் !!a{injection}. ஏனெனில்
$f(a, b) = f(c,d)$ எனில், எல்லா $n \in \Nat$-க்கும்
$a_n =
c_n$ மற்றும் $b_n = d_n$; ஆகவே $a = c$ மற்றும் $b =
d$.
\end{proof}

% SEGMENT OLP-0638-S03
ஆனால், டெடெகிண்ட் காண்டோருக்குச் சுட்டிக்காட்டியதுபோல்,
இது முதலில் எழுப்பிய கேள்விக்குப் பதிலளிக்கவில்லை.
$0.\dot{1}\dot{0} =
0.1010101010\ldots$ என்பதை எடுத்துக்கொள்ளுங்கள்.
$f(a,b) = 0.\dot{1}\dot{0}$ ஆக வேண்டுமெனில்:
\begin{align*}
a&= 0.\dot{1}\dot{1} = 0.111111\ldots\\
b&= 0
\end{align*}
ஆக வேண்டும். ஆனால் $a = 0.\dot{1}\dot{1} = 1$.
எனவே “$a$-ஐயும் $b$-ஐயும் இரும எண் குறியீட்டில்
எழுதுங்கள்” என்று சொல்லும்போது, \emph{எந்தக்}
குறியீட்டைத் தேர்ந்தெடுக்கிறோம் என்பதையும் முடிவுசெய்ய வேண்டும்.
$f$ ஒரு \emph{சார்பாக} இருக்க வேண்டுமென்றால்,
இரண்டு சாத்தியமான குறியீடுகளில் \emph{ஒன்றை மட்டுமே}
பயன்படுத்த முடியும். எடுத்துக்காட்டாக, எளிய குறியீட்டைத்
தேர்ந்து $a$-ஐ “$1.000\ldots$” என்று எழுதினால்,
$f(a, b) = 0.\dot{1}\dot{0}$ ஆகும் எந்த
$\tuple{a, b}$ இணையும் இல்லை.

% SOURCE CORRECTION TA-HIS-003: the displayed expansions and construction are binary, not decimal.
சுருக்கமாகச் சொன்னால், சில மீள்வரும் இரும விரிவாக்கங்கள்
சாத்தியம் என்பதால் காண்டோரின் $f$ சார்பு
!!a{injection} என்றாலும் !!a{surjection} \emph{அல்ல}
என்று டெடெகிண்ட் சுட்டிக்காட்டினார்.
எனவே காண்டோர் காட்டியது
$\cardle{\unitsquare}{\unitline}$ என்பதை மட்டுமே;
$\cardeq{\unitsquare}{\unitline}$ என்பதை \emph{அல்ல}.

% SEGMENT OLP-0638-S04
கிட்டத்தட்ட உடனே காண்டோர் டெடெகிண்டிற்குப் பதிலெழுதி,
நிறுவலைப் பின்வருமாறு நிறைவு செய்யலாம் என்று சாராம்சத்தில்
கூறினார்:

\begin{proof}[நிறுவல்: நிறைவு.]
$\cardle{\unitsquare}{\unitline}$ என்று காட்டிவிட்டோம்.
ஆனால் $\unitline$-இலிருந்து $\unitsquare$-க்கு
ஓர் !!a{injection} இருப்பது வெளிப்படை:
கோட்டைச் சதுரத்தின் ஒரு பக்கத்தின் மேல் வைத்தாலே போதும்.
ஆகவே $\cardle{\unitline}{\unitsquare}$ மற்றும்
$\cardle{\unitsquare}{\unitline}$. ஷ்ரோடர்--பெர்ன்ஸ்டைன்
தேற்றத்தின்படி
(\olref[sfr][siz][sb]{thm:schroder-bernstein}),
$\cardeq{\unitline}{\unitsquare}$.
\end{proof}

ஆனால் அந்தக் காலத்தில் ஷ்ரோடர்--பெர்ன்ஸ்டைன் தேற்றம்
இன்னும் நிறுவப்படாததால், காண்டோர் கடைசி வரியை
இவ்வாறு நிறைவு செய்ய இயலவில்லை. இதை ஒரு பொதுவான
ஊகமாக அவர் பின்னர் முன்வைத்தபோதும், 1897-இல்தான்
அது திருப்திகரமாக நிறுவப்பட்டது. (ஆகவே 1877-இலேயே
பின்னர் காண்டோர், ஷ்ரோடர்--பெர்ன்ஸ்டைன் தேற்றத்தைச்
சாராத \olref{thm:cantorplane} என்பதற்கான வேறொரு
நிறுவலை அளித்தார்.)

\end{document}
